%% chapters/04-metodologi.tex
%% Bab IV — Metodologi Penelitian
%% Sumber: writings/experiment-design.md (§1–§12)
%% Aturan: .cursor/rules/disertation/09-bab4-metodologi.mdc

\chapter{Metodologi Penelitian}
\label{bab:metodologi}

Bab ini menguraikan secara teknis dan \emph{reproducible} seluruh
desain metodologi penelitian, mencakup dataset yang digunakan, alur
praproses sinyal getaran, arsitektur model yang diusulkan, modul
interpretabilitas, protokol pelatihan \emph{multi-seed}, serta rancangan
eksperimen yang terstruktur.
Seluruh konfigurasi pada bab ini mencerminkan eksperimen aktual yang
dijalankan pada infrastruktur GPU cloud (NVIDIA A40, 46~GB VRAM).

\section{Kerangka Penelitian}
\label{sec:bab4_kerangka}

Penelitian ini disusun dalam dua pilar yang saling melengkapi.
\textbf{Pilar~1} berfokus pada performa prediksi sisa umur pakai
(\emph{remaining useful life}, RUL) menggunakan arsitektur hibrida
\emph{Mamba-xLSTM-Net}.
\textbf{Pilar~2} berfokus pada interpretabilitas representasi laten model
melalui \emph{Sparse Autoencoder} (SAE) dan pemetaan ke frekuensi
karakteristik fisik bantalan.
Kerangka penelitian secara keseluruhan ditunjukkan pada
Gambar~\ref{fig:bab4_kerangka}.

\begin{figure}[ht]
\centering
\resizebox{\textwidth}{!}{%
\begin{tikzpicture}[
    node distance = 0.5cm and 1.2cm,
    proc/.style  = {rectangle, rounded corners=3pt, draw, thick, fill=white,
                    align=center, minimum height=0.85cm, minimum width=2.8cm,
                    font=\small},
    data/.style  = {proc, fill=gray!15},
    pilar1s/.style = {proc, fill=blue!10},
    pilar2s/.style = {proc, fill=green!10},
    metric/.style  = {proc, fill=yellow!15},
    arr/.style  = {-{Stealth[length=5pt]}, thick},
  ]
  \node[data]   (sensor) {Akselerometer\\{\small PHM2012, XJTU-SY}};
  \node[proc, right=of sensor]  (feat)   {Ekstraksi Fitur HI\\{\scriptsize band-energy + statistik}};
  \node[proc, right=of feat]    (window) {Jendela Sekuens\\{\scriptsize $L = 32$}};

  \draw[arr] (sensor) -- (feat);
  \draw[arr] (feat)   -- (window);

  \node[pilar1s, below=1.0cm of window] (mxnet) {Mamba-xLSTM-Net\\{\scriptsize Pilar~1}};
  \node[metric,  right=of mxnet] (rul) {Prediksi RUL $\hat{r}_t$\\{\scriptsize RMSE, MAE, $R^2$}};

  \draw[arr] (window.south) -- (mxnet.north);
  \draw[arr] (mxnet) -- (rul);

  \node[pilar2s, below=1.0cm of mxnet] (sae) {Top-$k$ SAE\\{\scriptsize Pilar~2 post-hoc}};
  \node[metric,  right=of sae] (bpfx) {Pemetaan BPFx\\{\scriptsize hit-rate, korelasi Pearson}};

  \draw[arr, dashed] (mxnet.south) -- (sae.north)
        node[midway, right, font=\scriptsize]{laten $\mathbf{z}$};
  \draw[arr] (sae) -- (bpfx);

  \node[font=\scriptsize, text=blue!60!black, left=0.1cm of mxnet] {$\uparrow$ P1};
  \node[font=\scriptsize, text=green!50!black, left=0.1cm of sae]  {$\uparrow$ P2};
\end{tikzpicture}%
}
\caption{Kerangka penelitian keseluruhan: sinyal akselerometer diproses menjadi
         fitur \emph{health indicator} (HI) dan sekuens berjendela, kemudian
         dimasukkan ke Mamba-xLSTM-Net untuk prediksi RUL (Pilar~1); representasi
         laten model secara \emph{post-hoc} diinterpretasi oleh Top-$k$ SAE dan
         dipetakan ke frekuensi karakteristik fisik bantalan (Pilar~2).}
\label{fig:bab4_kerangka}
\end{figure}

\section{Dataset}
\label{sec:bab4_dataset}

Eksperimen menggunakan dua dataset publik \emph{run-to-failure} yang
paling banyak dipakai dalam literatur prognostik bantalan.
Penggunaan dua dataset sekaligus adalah keputusan disengaja: PHM2012
menyediakan jumlah bantalan lebih besar dengan rekaman pendek dan padat,
sementara XJTU-SY menyediakan rekaman panjang dengan label mode kegagalan
eksplisit.
Kombinasi keduanya memberikan basis empiris yang lebih kuat dibandingkan
menggunakan satu dataset saja.

\subsection{PHM2012 (FEMTO-PRONOSTIA)}
\label{subsec:bab4_phm2012}

Dataset PHM2012 merupakan hasil kompetisi \emph{Prognostics and Health
Management} IEEE 2012 yang dikembangkan oleh FEMTO-ST Institute,
Université de Franche-Comté, menggunakan \emph{test rig} PRONOSTIA
\citetitb{Nectoux2012PRONOSTIA}.
Rincian atribut dataset ditampilkan pada Tabel~\ref{tab:bab4_phm2012}.

\begin{table}[ht]
\centering
\caption{Karakteristik dataset PHM2012 (FEMTO-PRONOSTIA).}
\label{tab:bab4_phm2012}
\small
\begin{tabular}{@{}ll@{}}
\toprule
\textbf{Atribut} & \textbf{Nilai} \\
\midrule
Sumber                  & FEMTO-ST Institute, Université de Franche-Comté \\
\emph{Test rig}         & PRONOSTIA \\
Bantalan uji            & NSK 6804 (\emph{deep groove ball bearing}) \\
Kondisi operasi         & 3 (1800~rpm/4000~N; 1650~rpm/4200~N; 1500~rpm/5000~N) \\
Total bantalan          & 17 (6 \emph{training set}, 11 \emph{test set}) \\
Sensor                  & Akselerometer horizontal + vertikal (25{,}6~kHz), termokopel \\
Durasi rekaman          & 0{,}1 detik per rekaman; interval 10 detik \\
Sampel per rekaman      & 2.560 \\
Skema label RUL         & Linear dari 1{,}0 (sehat) ke 0{,}0 (EOL, \emph{threshold} 20~g) \\
\bottomrule
\end{tabular}
\end{table}

Kompetisi PHM2012 menggunakan \emph{scoring function} asimetris yang
menghukum prediksi terlalu dini lebih berat daripada prediksi terlambat
\citetitb{Nectoux2012PRONOSTIA}.
Konvensi pembagian data kompetisi ini dipertahankan dalam penelitian ini
agar hasil dapat dibandingkan langsung dengan literatur \emph{benchmark}
yang telah dipublikasikan.

Gambar~\ref{fig:bab4_dataset_lifecycle} mengilustrasikan profil nilai RMS
akselerometer sebuah bantalan PHM2012 (Bearing~1\_1) sepanjang siklus
hidupnya, memperlihatkan fase sehat yang panjang diikuti fase degradasi
akseleratif menjelang kegagalan. Profil ini secara visual mengonfirmasi
asumsi \emph{piecewise} RUL (Persamaan~\ref{eq:bab3_label_piecewise}) yang
digunakan dalam penelitian ini, di mana penurunan RUL baru mulai pada titik
\emph{degradation onset} $t_d$ yang terdeteksi otomatis.

\begin{figure}[ht]
  \centering
  % TODO(produce-figure): scripts/plot_hi_lifecycle.py
  % Data: data-bearing/ieee-phm-2012/Learning_set/Bearing1_1/acc_*.csv
  % Plot RMS rolling window vs waktu, annotate t_d dan EOL threshold (20 g).
  \fbox{\parbox{0.88\textwidth}{\centering\vspace{1.8cm}
    \textit{Gambar IV.\ref{fig:bab4_dataset_lifecycle} --- Siklus Hidup HI
    Bearing 1\_1 (PHM2012)}\\[6pt]
    \textit{Plot RMS akselerometer vs. rekaman ke-$t$, dengan anotasi titik
    \emph{degradation onset} $t_d$ (garis vertikal putus-putus) dan ambang
    kegagalan 20~g (garis horizontal merah).}\\[4pt]
    \textit{Akan dirender oleh \texttt{scripts/plot\_hi\_lifecycle.py}
    sebelum sidang.}
    \vspace{1.8cm}}}
  \caption{Profil \emph{health indicator} (nilai RMS akselerometer) sepanjang
           siklus hidup Bearing~1\_1 dari dataset PHM2012. Fase sehat panjang
           ($r_t = 1$) diikuti fase degradasi akseleratif yang relatif singkat;
           titik \emph{degradation onset} $t_d$ ditentukan otomatis melalui
           metode $3\sigma$ geser, dan kegagalan akhir (\emph{end of life})
           ditetapkan saat RMS $\geq 20$~g.}
  \label{fig:bab4_dataset_lifecycle}
\end{figure}

\subsection{XJTU-SY}
\label{subsec:bab4_xjtusy}

Dataset XJTU-SY dikembangkan oleh Xi'an Jiaotong University bersama
Sumyoung Technology menggunakan \emph{test rig} akselerasi umur pakai
bantalan gelinding \citetitb{Wang2020XJTU}.
Rincian atribut dataset ditampilkan pada Tabel~\ref{tab:bab4_xjtusy}.

\begin{table}[ht]
\centering
\caption{Karakteristik dataset XJTU-SY.}
\label{tab:bab4_xjtusy}
\small
\begin{tabular}{@{}ll@{}}
\toprule
\textbf{Atribut} & \textbf{Nilai} \\
\midrule
Sumber              & Xi'an Jiaotong University \& Sumyoung Technology \\
Bantalan uji        & LDK UER204 (\emph{deep groove ball bearing}) \\
Kondisi operasi     & 3 (2100~rpm/12~kN; 2250~rpm/11~kN; 2400~rpm/10~kN) \\
Total bantalan      & 15 (kondisi 1: 5; kondisi 2: 5; kondisi 3: 5) \\
Sensor              & Akselerometer horizontal + vertikal (25{,}6~kHz) \\
Durasi rekaman      & 1{,}28 detik per rekaman; interval 1 menit \\
Sampel per rekaman  & 32.768 \\
Label mode kegagalan & \emph{Outer race}, \emph{inner race}, \emph{cage} (per bantalan) \\
Catatan             & Kondisi 3 tidak tersedia; hanya kondisi 1 dan 2 digunakan \\
\bottomrule
\end{tabular}
\end{table}

Ketersediaan label mode kegagalan per bantalan pada XJTU-SY menjadi aset
penting untuk validasi interpretabilitas pada Pilar~2: bila SAE memetakan
fitur ke BPFO atau BPFI, label mode kegagalan eksplisit memungkinkan
uji silang apakah pemetaan tersebut konsisten dengan jenis cacat aktual
yang didokumentasikan pemilik dataset.

\section{Praproses Sinyal dan Ekstraksi Fitur}
\label{sec:bab4_praproses}

\emph{Pipeline} praproses terdiri dari tiga tahap utama yang dirancang
agar sinyal getaran mentah dapat dikonsumsi oleh model \emph{deep learning}
dalam bentuk vektor fitur \emph{Health Indicator} (HI) yang ringan secara
komputasi tanpa kehilangan informasi diagnostik.

\subsection{Dekomposisi Energi Pita Frekuensi}
\label{subsec:bab4_bandenergi}

Setiap rekaman getaran mentah dilewatkan melalui \emph{filter bank} yang
membaginya ke dalam $N_{\text{bands}} = 5$ sub-pita frekuensi.
Untuk setiap sub-pita, dihitung energi RMS (\emph{root-mean-square})
sebagai representasi distribusi energi spektral:

\begin{equation}\label{eq:bab4_band_energy}
E_{k} = \sqrt{\frac{1}{N}\sum_{i=1}^{N} x_{k,i}^{2}}, \quad k = 1, 2, \ldots, N_{\text{bands}},
\end{equation}

\noindent di mana $x_{k,i}$ adalah sampel sinyal pada sub-pita ke-$k$
dan $N$ adalah jumlah sampel per sub-pita.
Pendekatan ini meminjam dari teori \emph{envelope analysis} klasik yang
menunjukkan bahwa cacat bantalan menghasilkan modulasi amplitudo yang
terkonsentrasi pada pita frekuensi tertentu, umumnya di sekitar resonansi
tinggi \citetitb{Lei2018Machinery}.

\subsection{Skema Label RUL}
\label{subsec:bab4_label}

Label HI \emph{ground truth} dihasilkan dengan skema linear dari 1{,}0
(saat rekaman pertama) ke 0{,}0 (saat \emph{end-of-life}), ternormalisasi
terhadap total panjang hidup tiap bantalan:

\begin{equation}\label{eq:bab4_rul_linear}
r_{t} = \frac{T - t}{T}, \quad t = 0, 1, \ldots, T-1,
\end{equation}

\noindent di mana $T$ adalah total jumlah rekaman dari rekaman pertama
hingga rekaman terakhir sebelum kegagalan, dan $t$ adalah indeks rekaman.
Untuk meredam \emph{noise} pada fase awal, digunakan \emph{exponential
moving average} dengan koefisien pemulusan $\alpha_{\text{smooth}} = 0{,}10$.

\subsection{Windowing Temporal}
\label{subsec:bab4_window}

Model menerima input berupa \emph{sliding window} yang berisi
$L_{\text{win}}$ rekaman berurutan.
Ukuran \emph{window} dan \emph{stride} berbeda per dataset untuk
menyeimbangkan kepadatan informasi dan biaya komputasi, sebagaimana
ditampilkan pada Tabel~\ref{tab:bab4_window}.

\begin{table}[ht]
\centering
\caption{Parameter \emph{windowing} per dataset.}
\label{tab:bab4_window}
\small
\begin{tabular}{@{}lSS@{}}
\toprule
\textbf{Parameter} & \textbf{PHM2012} & \textbf{XJTU-SY} \\
\midrule
$L_{\text{win}}$ (rekaman)  & 64    & 32    \\
$s_{\text{train}}$          & 1     & 1     \\
$s_{\text{eval}}$           & 32    & 1     \\
\bottomrule
\end{tabular}
\end{table}

\emph{Stride} pelatihan $s_{\text{train}} = 1$ menjamin augmentasi
alami melalui \emph{overlapping windows}.
\emph{Stride} evaluasi pada PHM2012 $s_{\text{eval}} = 32$ dipilih agar
evaluasi konservatif (\emph{non-overlapping}), namun pada XJTU-SY
$s_{\text{eval}} = 1$ karena ukuran \emph{test set} yang kecil
memerlukan setiap \emph{window} untuk diestimasi.

\section{Pembagian Data}
\label{sec:bab4_split}

Pembagian data dilakukan secara berbasis bantalan (\emph{bearing-wise})
untuk mencegah kebocoran data (\emph{data leakage}).
Sampel dari satu bantalan yang sama tidak diizinkan tersebar ke set
pelatihan dan pengujian, karena karakteristik degradasi unik tiap
bantalan dapat dipelajari secara dangkal oleh model.
Statistik normalisasi (mean, deviasi standar) dihitung \emph{hanya}
pada \emph{training set} dan diterapkan ke \emph{validation} serta
\emph{test set} tanpa modifikasi lebih lanjut.

Gambar~\ref{fig:bab4_split_diagram} memperlihatkan skema pembagian
\emph{bearing-wise} untuk kedua dataset secara berdampingan. Prinsip
utama yang dijaga adalah bahwa identitas bantalan tidak pernah
bersilang antar-set, sehingga metrik evaluasi pada \emph{test set}
mencerminkan kemampuan generalisasi model terhadap bantalan yang
belum pernah dilihat selama pelatihan.

\begin{figure}[ht]
  \centering
  \resizebox{\textwidth}{!}{%
  \begin{tikzpicture}[
      node distance = 0.35cm and 0.5cm,
      bear/.style = {rectangle, draw, thick, text centered, minimum height=0.6cm,
                     minimum width=1.1cm, font=\scriptsize},
      train/.style = {bear, fill=blue!15},
      val/.style   = {bear, fill=orange!20},
      test/.style  = {bear, fill=red!12},
      hdr/.style   = {font=\small\bfseries, text centered},
      lbl/.style   = {font=\scriptsize, anchor=west},
    ]
    \node[hdr] (ph) {PHM2012};
    \node[train, below=0.3cm of ph] (p1) {1\_1};
    \node[train, right=of p1] (p2) {1\_2};
    \node[train, right=of p2] (p3) {1\_4};
    \node[train, right=of p3] (p4) {2\_1};
    \node[train, right=of p4] (p5) {2\_3};
    \node[train, right=of p5] (p6) {2\_5};
    \node[train, right=of p6] (p7) {3\_1};
    \node[val,   right=0.6cm of p7] (pv1) {1\_5};
    \node[val,   right=of pv1] (pv2) {2\_2};
    \node[test,  right=0.6cm of pv2] (pt1) {1\_3};
    \node[test,  right=of pt1] (pt2) {$\cdots$};
    \node[test,  right=of pt2] (pt3) {3\_3};

    % Row labels
    \node[lbl] at ($(p1.west)+(-0.1cm,0)$) {};
    % Legend
    \node[train, below=0.7cm of p1, minimum width=0.7cm] (lg1) {};
    \node[right=0.1cm of lg1, font=\scriptsize] {Latih (7)};
    \node[val,   right=0.8cm of lg1, minimum width=0.7cm] (lg2) {};
    \node[right=0.1cm of lg2, font=\scriptsize] {Validasi (2)};
    \node[test,  right=0.8cm of lg2, minimum width=0.7cm] (lg3) {};
    \node[right=0.1cm of lg3, font=\scriptsize] {Uji (8)};

    % --- XJTU-SY ---
    \node[hdr, below=2.0cm of ph] (xj) {XJTU-SY};
    \node[train, below=0.3cm of xj] (x1) {1\_1};
    \node[train, right=of x1] (x2) {1\_2};
    \node[train, right=of x2] (x3) {1\_3};
    \node[train, right=of x3] (x4) {2\_1};
    \node[train, right=of x4] (x5) {2\_2};
    \node[val,   right=0.6cm of x5] (xv) {1\_4};
    \node[test,  right=0.6cm of xv] (xt1) {1\_5};
    \node[test,  right=of xt1] (xt2) {2\_3};
  \end{tikzpicture}%
  }
  \caption{Skema pembagian \emph{bearing-wise} pada PHM2012 (atas) dan
           XJTU-SY (bawah). Setiap kotak mewakili satu bantalan; warna
           menunjukkan set: biru = latih, oranye = validasi, merah = uji.
           Tidak ada bantalan yang bersilang antar-set, memastikan evaluasi
           bebas kebocoran data.}
  \label{fig:bab4_split_diagram}
\end{figure}

\subsection{Pembagian PHM2012}
\label{subsec:bab4_split_phm}

Pembagian \emph{bearing-wise} yang diterapkan pada dataset PHM2012 dirinci
pada Tabel~\ref{tab:bab4_split_phm}. Pemilihan bantalan untuk masing-masing
set mengikuti konvensi kompetisi PHM2012 sehingga memungkinkan perbandingan
langsung dengan literatur.

\begin{table}[ht]
\centering
\caption{Pembagian set data PHM2012 (bearing-wise).}
\label{tab:bab4_split_phm}
\small
\begin{tabular}{@{}lll@{}}
\toprule
\textbf{Set}       & \textbf{Bantalan}                              & \textbf{Jumlah} \\
\midrule
\emph{Training}    & 1\_1, 1\_2, 1\_4, 2\_1, 2\_3, 2\_5, 3\_1      & 7 bantalan \\
\emph{Validation}  & 1\_5, 2\_2                                     & 2 bantalan \\
\emph{Test}        & 1\_3, 1\_6, 1\_7, 2\_4, 2\_6, 2\_7, 3\_2, 3\_3 & 8 bantalan \\
\bottomrule
\end{tabular}
\end{table}

Pembagian ini mengikuti konvensi eksperimen awal di literatur PHM2012
sehingga hasil dapat dibandingkan langsung dengan tabel \emph{benchmark}
yang dipublikasikan \citetitb{Nectoux2012PRONOSTIA}.

\subsection{Pembagian XJTU-SY}
\label{subsec:bab4_split_xjtu}

Pembagian \emph{bearing-wise} untuk XJTU-SY ditampilkan pada
Tabel~\ref{tab:bab4_split_xjtu}. Karena kondisi 3 tidak tersedia
pada infrastruktur yang digunakan, hanya 10 bantalan dari kondisi 1
dan 2 yang diikutsertakan.

\begin{table}[ht]
\centering
\caption{Pembagian set data XJTU-SY (bearing-wise, 2 kondisi aktif).}
\label{tab:bab4_split_xjtu}
\small
\begin{tabular}{@{}lll@{}}
\toprule
\textbf{Set}       & \textbf{Bantalan}           & \textbf{Jumlah} \\
\midrule
\emph{Training}    & 1\_1, 1\_2, 1\_3, 2\_1, 2\_2 & 5 bantalan \\
\emph{Validation}  & 1\_4                          & 1 bantalan \\
\emph{Test}        & 1\_5, 2\_3                    & 2 bantalan \\
\bottomrule
\end{tabular}
\end{table}

\emph{Test set} XJTU-SY hanya berisi dua bantalan karena keterbatasan
ketersediaan data (kondisi 3 tidak dapat diakses).
Konsekuensinya, variansi metrik antar-\emph{seed} pada XJTU-SY relatif
tinggi dan R\textsuperscript{2} sering negatif karena \emph{sum of
squares total} yang kecil.
RMSE tetap dipakai sebagai metrik utama untuk mempertahankan
komparabilitas dengan literatur.

\section{Arsitektur Mamba-xLSTM-Net}
\label{sec:bab4_arsitektur}

\emph{Mamba-xLSTM-Net} merupakan arsitektur fokus disertasi yang
mengintegrasikan dua blok pemrosesan sekuens: (1) \emph{Selective
State-Space Model} (Mamba-3) yang memproses konteks temporal jangka
panjang secara efisien, dan (2) \emph{Matrix-Memory Long Short-Term
Memory} (mLSTM) yang mempertahankan presisi rekuren pada perubahan
tajam menjelang kegagalan.

\subsection{Justifikasi Desain Hibrida}
\label{subsec:bab4_justifikasi}

Pertanyaan mendasar dalam desain arsitektur prediksi RUL adalah: \emph{mengapa
arsitektur hibrida, bukan arsitektur tunggal yang semakin diperdalam?} Tinjauan
kritis pada Bab~\ref{bab:tinjauan-pustaka}
menunjukkan bahwa tidak ada satu keluarga arsitektur yang secara tunggal mengatasi
semua tantangan yang relevan untuk sinyal degradasi bantalan. Bagian ini
memformalisasi justifikasi tersebut dalam kerangka komplementaritas dua komponen.

Sinyal degradasi bantalan menunjukkan dua rezim temporal yang memiliki karakteristik
yang kontras secara fundamental. Pada \textbf{fase degradasi gradual} (awal hingga
sekitar $70\%$ siklus hidup), amplitudo getaran meningkat perlahan, korelasi
temporal mencakup jangkauan panjang, dan pola global lebih informatif daripada
perubahan lokal. Pada \textbf{fase akselerasi degradasi} (sekitar $30\%$ akhir
siklus hidup), karakter sinyal berubah secara fundamental: terjadi lonjakan
amplitudo yang dramatis, fitur kurtosis dan faktor puncak meningkat tajam, dan
model harus mampu merespons dengan revisi besar terhadap estimasi RUL dalam
waktu singkat. Tidak ada komponen tunggal yang optimal untuk kedua rezim
sekaligus.

Tabel~\ref{tab:bab4_komplementer} menyajikan perbandingan lima aspek kritis
antara BiMamba dan xLSTM, menunjukkan bagaimana kelemahan satu komponen
dilengkapi oleh kekuatan komponen lainnya, dan bagaimana sinergi dalam
Mamba-xLSTM-Net mengeksploitasi komplementaritas ini.

\begin{table}[ht]
  \centering
  \caption{Perbandingan komplementer BiMamba dan xLSTM (sLSTM + mLSTM) pada
           lima aspek kritis, dan sinergi yang dihasilkan oleh Mamba-xLSTM-Net.}
  \label{tab:bab4_komplementer}
  \footnotesize
  \begin{tabularx}{\textwidth}{@{}>{\raggedright\arraybackslash}p{2.4cm} X X X@{}}
    \toprule
    \textbf{Aspek} & \textbf{BiMamba} & \textbf{xLSTM (sLSTM+mLSTM)} &
    \textbf{Sinergi Mamba-xLSTM-Net} \\
    \midrule
    Kompleksitas komputasi
      & $\mathcal{O}(L \cdot d_{\text{state}})$ — linear terhadap panjang sekuens
      & $\mathcal{O}(L \cdot d^2)$ untuk mLSTM — kuadratik terhadap dimensi
      & BiMamba menangani konteks panjang; xLSTM beroperasi pada jendela lokal
        yang lebih pendek \\[6pt]
    Sensitivitas akhir umur
      & Rendah — \emph{selective scan} yang smooth kurang responsif terhadap
        lonjakan mendadak
      & Tinggi — \emph{exponential gating} tanpa saturasi merespons lonjakan
        amplitudo secara agresif \citetitb{Beck2024xLSTM}
      & Gate fusi mengalihkan bobot ke xLSTM ketika estimasi RUL $\hat{r}_t < 0{,}3$
        \\[6pt]
    Memori asosiatif
      & $\mathcal{O}(d_{\text{state}})$ — state tersembunyi vektor $N$-dimensi
      & $\mathcal{O}(d^2)$ — memori matriks mLSTM menyimpan multi-asosiasi
        kunci-nilai
      & mLSTM menyimpan korelasi fitur-BPFx; BiMamba menyimpan tren global
        \\[6pt]
    Bidireksionalitas
      & Native — dua \emph{scan} paralel (maju + mundur) memberikan konteks
        penuh jendela tanpa kausalitas paksa
      & Tidak native — sLSTM dan mLSTM rekurens kausal
      & BiMamba memberikan konteks dua arah untuk seluruh jaringan;
        xLSTM memanfaatkan representasi yang sudah diperkaya \\[6pt]
    Beban komputasi pelatihan
      & Rendah — \emph{parallel scan} yang sangat efisien pada GPU,
        $2\times$ overhead untuk dua arah
      & Moderat — kepala (\emph{heads}) paralel, tetapi sekuensial dalam arah
        rekurens
      & Kedua cabang dilatih secara bersamaan; total parameter $<\!1$M
        (lihat Tabel~\ref{tab:bab4_hyperp_mamba}) \\
    \bottomrule
  \end{tabularx}
\end{table}

Justifikasi tambahan bersumber dari analisis mekanistik pada Bab~\ref{bab:dasar-teori}:
limitasi saturasi sigmoid LSTM yang menjadi motivasi \emph{exponential gating}
xLSTM diuraikan pada subbab~\ref{subsec:bab3_lstm}, keterbatasan kapasitas
vektor memori yang memotivasi memori matriks mLSTM diuraikan pada
subbab~\ref{subsec:bab3_slstm}, dan argumen untuk konfigurasi bidireksional
Mamba diuraikan pada subbab~\ref{subsec:bab3_bimamba}. Tinjauan literatur pada Bab~\ref{bab:tinjauan-pustaka} mengkonfirmasi bahwa tidak ada satu pun
keluarga arsitektur tunggal yang menutup semua dimensi kelemahan.
Gambar~\ref{fig:bab4_rasional} mengilustrasikan secara visual mekanisme
\emph{gated fusion} yang mengintegrasikan dua cabang paralel arsitektur.

\begin{figure}[ht]
  \centering
  \resizebox{\textwidth}{!}{%
  \begin{tikzpicture}[
      node distance = 0.6cm and 1.6cm,
      inp/.style  = {rectangle, rounded corners=3pt, draw, thick, fill=gray!15,
                     align=center, minimum height=0.8cm, minimum width=2.4cm,
                     font=\small},
      bimamba/.style = {inp, fill=blue!12},
      xlstm/.style   = {inp, fill=green!12},
      fusion/.style  = {inp, fill=orange!15, minimum width=2.6cm},
      outbox/.style  = {inp, fill=yellow!20},
      arr/.style = {-{Stealth[length=5pt]}, thick},
    ]
    \node[inp] (x) {$\mathbf{x}_{1:L}$ (HI sekuens)};

    % Two parallel branches
    \node[bimamba, above right=0.9cm and 1.8cm of x] (bm)
         {BiMamba\\{\scriptsize $\times 4$ blok, konteks panjang}};
    \node[xlstm, below right=0.9cm and 1.8cm of x] (xl)
         {xLSTM (sLSTM + mLSTM)\\{\scriptsize $\times 4$ blok, dinamika lokal}};

    \draw[arr] (x.east) -- (bm.west);
    \draw[arr] (x.east) -- (xl.west);

    % Fusion
    \coordinate (bm_xl_mid) at ($(bm)!0.5!(xl)$);
    \node[fusion, right=1.8cm of bm_xl_mid] (fus)
         {\emph{Gated Fusion}\\{\scriptsize $\alpha_t \in [0,1]$}};

    \draw[arr] (bm.east) -- (fus.north west);
    \draw[arr] (xl.east) -- (fus.south west);

    % Output
    \node[outbox, right=of fus] (rul) {$\hat{r}_t$ (RUL)};
    \draw[arr] (fus) -- (rul);

    % Annotations
    \node[font=\scriptsize, text=blue!60!black, above=0.1cm of bm]
         {tren jangka panjang};
    \node[font=\scriptsize, text=green!50!black, below=0.1cm of xl]
         {lonjakan akhir umur};
  \end{tikzpicture}%
  }
  \caption{Rasional desain hibrida Mamba-xLSTM-Net: dua cabang paralel
           memproses sekuens HI secara bersamaan. Cabang BiMamba ($\times 4$ blok)
           menangkap tren degradasi jangka panjang dan konteks global, sedangkan
           cabang xLSTM ($\times 4$ blok) mempertahankan sensitivitas dinamika
           lokal dan lonjakan akhir umur. Keduanya digabungkan oleh mekanisme
           \emph{gated fusion} adaptif dengan bobot $\alpha_t$ yang bergantung
           pada input untuk menghasilkan prediksi RUL $\hat{r}_t$.}
  \label{fig:bab4_rasional}
\end{figure}

\subsection{Detail Arsitektur}
\label{subsec:bab4_detail}

Input berupa tensor berdimensi $B \times L \times F$ (batch × panjang
\emph{window} × jumlah fitur) diproses melalui lapisan proyeksi
konvolusional (\emph{convolutional front-end}) sebelum masuk ke tumpukan
blok Mamba.
Keluaran blok Mamba kemudian diproses oleh tumpukan blok mLSTM sebelum
akhirnya direduksi melalui \emph{prediction head} (satu lapisan linear)
menjadi prediksi RUL skalar. Gambar~\ref{fig:bab4_arsitektur} menunjukkan
diagram blok arsitektur lengkap dari lapisan proyeksi hingga keluaran prediksi.

\begin{figure}[ht]
  \centering
  \resizebox{0.95\textwidth}{!}{%
  \begin{tikzpicture}[
      node distance = 0.45cm and 0.0cm,
      blk/.style  = {rectangle, rounded corners=3pt, draw, thick, fill=white,
                     align=center, minimum height=0.85cm, minimum width=3.2cm,
                     font=\small},
      bm/.style   = {blk, fill=blue!10},
      xl/.style   = {blk, fill=green!10},
      fus/.style  = {blk, fill=orange!15},
      head/.style = {blk, fill=yellow!15},
      io/.style   = {blk, fill=gray!15},
      arr/.style  = {-{Stealth[length=5pt]}, thick},
    ]
    \node[io]  (inp)   {Input $B \times L \times F$};
    \node[blk, below=of inp]   (proj)  {\emph{Conv Front-end}\\{\scriptsize proj ke $d_{\rm model}=128$}};
    \draw[arr] (inp)  -- (proj);

    % BiMamba stack
    \node[bm, below=0.6cm of proj]  (bm1) {BiMamba Blok 1};
    \node[bm, below=of bm1] (bm2) {BiMamba Blok 2};
    \node[bm, below=of bm2] (bm3) {BiMamba Blok 3};
    \node[bm, below=of bm3] (bm4) {BiMamba Blok 4};
    \draw[arr] (proj) -- (bm1);
    \draw[arr] (bm1) -- (bm2);
    \draw[arr] (bm2) -- (bm3);
    \draw[arr] (bm3) -- (bm4);

    % xLSTM stack
    \node[xl, right=2.0cm of bm1] (xl1) {mLSTM Blok 1};
    \node[xl, below=of xl1] (xl2) {mLSTM Blok 2};
    \node[xl, below=of xl2] (xl3) {mLSTM Blok 3};
    \node[xl, below=of xl3] (xl4) {mLSTM Blok 4};
    \draw[arr] (proj.east) -- ++(0.5cm,0) |- (xl1.west);
    \draw[arr] (xl1) -- (xl2);
    \draw[arr] (xl2) -- (xl3);
    \draw[arr] (xl3) -- (xl4);

    % Fusion
    \coordinate (bm4_xl4_mid) at ($(bm4)!0.5!(xl4)$);
    \node[fus, below=0.9cm of bm4_xl4_mid] (fu) {\emph{Gated Fusion} ($\alpha_t$)};
    \draw[arr] (bm4.south) -- (fu.north west);
    \draw[arr] (xl4.south) -- (fu.north east);

    % Head
    \node[head, below=0.5cm of fu] (hd) {\emph{Regression Head} (Linear)};
    \draw[arr] (fu) -- (hd);

    \node[io, below=0.4cm of hd] (out) {Output $\hat{r}_t \in [0,1]$};
    \draw[arr] (hd) -- (out);

    % Annotations
    \node[font=\scriptsize, text=blue!60!black, left=0.1cm of bm1] {$\leftarrow$ \footnotesize konteks panjang};
    \node[font=\scriptsize, text=green!50!black, right=0.1cm of xl1] {\footnotesize dinamika lokal $\rightarrow$};
  \end{tikzpicture}%
  }
  \caption{Diagram blok arsitektur Mamba-xLSTM-Net. Input sekuens HI
           diproyeksikan oleh lapisan konvolusional, kemudian diproses secara
           paralel oleh empat blok BiMamba (cabang kiri, konteks panjang) dan
           empat blok mLSTM (cabang kanan, dinamika lokal). Keluaran kedua cabang
           digabungkan oleh \emph{gated fusion} adaptif sebelum diregresikan
           oleh \emph{prediction head} menjadi RUL ternormalisasi.}
  \label{fig:bab4_arsitektur}
\end{figure}

Hiperparameter arsitektur disajikan pada Tabel~\ref{tab:bab4_hyperp_mamba}.

\begin{table}[ht]
\centering
\caption{Hiperparameter Mamba-xLSTM-Net.}
\label{tab:bab4_hyperp_mamba}
\small
\begin{tabular}{@{}ll@{}}
\toprule
\textbf{Hiperparameter} & \textbf{Nilai} \\
\midrule
\texttt{d\_model}          & 128    \\
\texttt{mamba\_d\_state}   & 128    \\
\texttt{mamba\_expand}     & 2 ($d_{\text{inner}} = 256$)    \\
Jumlah blok Mamba          & 4      \\
Jumlah \emph{head} mLSTM   & 4      \\
\texttt{dropout}           & 0{,}10 \\
\emph{Input window}        & 64 (PHM) / 32 (XJTU) \\
\emph{Input features}      & 5 \\
Total parameter            & 898.481 \\
\bottomrule
\end{tabular}
\end{table}

\section{Model Pembanding}
\label{sec:bab4_pembanding}

Untuk mengevaluasi keunggulan komparatif Mamba-xLSTM-Net, dua arsitektur
\emph{Tier-S} lainnya digunakan sebagai pembanding.
Pemilihan model pembanding dilakukan melalui fase \emph{pre-screening}
30~epoch pada 12 kandidat model; tiga model dengan kinerja terburuk
(PhaseMoE-xLSTM, Physics-N-BEATS, LiquidWave-RUL) dikecualikan dari
eksperimen utama.

\subsection{N-BEATS-xLSTM-RUL}
\label{subsec:bab4_nbeats}

N-BEATS-xLSTM-RUL menggabungkan blok-blok basis polinomial
\emph{physics-inspired} dengan \emph{xLSTM temporal front} dan
\emph{monotone head} untuk dekomposisi RUL yang dapat diinterpretasi
per komponen (tren, keausan, kejutan).
Hiperparameter: \texttt{hidden\_dim} = 96, 2 blok \emph{trend}, 2 blok
\emph{wear}, 2 blok \emph{shock}, derajat polinomial 4.
Total parameter: 459.592 (PHM) / 457.672 (XJTU).

\subsection{SparseGate-TCN-RUL}
\label{subsec:bab4_sparsegate}

SparseGate-TCN-RUL menggunakan TCN (\emph{Temporal Convolutional
Network}) dengan \emph{sparse gating} yang memaksa hanya sebagian kecil
\emph{channel} aktif per sampel, memberikan interpretabilitas
\emph{built-in} tanpa SAE \emph{post-hoc}.
Hiperparameter: \texttt{tcn\_\allowbreak channels} = [64, 64, 128, 128],
\texttt{lambda\_\allowbreak sparse} = 0{,}001,
\texttt{lambda\_\allowbreak entropy} = 0{,}001.
Total parameter: 249.192.

\section{Modul Interpretabilitas}
\label{sec:bab4_interpretabilitas}

Penelitian ini menggunakan tiga pendekatan interpretabilitas yang
saling melengkapi: \emph{Top-$k$ Sparse Autoencoder} (pilar utama),
SHAP (\emph{SHapley Additive exPlanations}), dan \emph{Integrated
Gradients} (IG).

\subsection{Top-$k$ Sparse Autoencoder}
\label{subsec:bab4_sae}

SAE dilatih secara \emph{post-hoc} pada \emph{hidden state} dari
\emph{checkpoint} terbaik Mamba-xLSTM-Net.
Arsitektur SAE mengikuti \emph{recipe} Anthropic \citetitb{Bricken2023Monosemanticity}
dengan \emph{expansion factor} $e = 8$:

\begin{equation}\label{eq:bab4_sae_latent}
\mathbf{f} = \text{Top}_{k}\!\left(\mathbf{W}_{\text{enc}} \mathbf{h} + \mathbf{b}_{\text{enc}}\right),
\end{equation}

\begin{equation}\label{eq:bab4_sae_recon}
\hat{\mathbf{h}} = \mathbf{W}_{\text{dec}}^{\top} \mathbf{f} + \mathbf{b}_{\text{dec}},
\end{equation}

\noindent di mana $\mathbf{h} \in \mathbb{R}^{128}$ adalah \emph{hidden
state} Mamba-xLSTM-Net, $\mathbf{f} \in \mathbb{R}^{1024}$ adalah vektor
fitur laten \emph{sparse} dengan tepat $k = 51$ elemen aktif
($\approx 5\%$ sparsity), dan $\hat{\mathbf{h}}$ adalah rekonstruksi.
Fungsi \emph{loss} SAE:

\begin{equation}\label{eq:bab4_sae_loss}
\mathcal{L}_{\text{SAE}} = \left\|\mathbf{h} - \hat{\mathbf{h}}\right\|_{2}^{2},
\end{equation}

\noindent tanpa regularisasi $\ell_1$ tambahan karena sparsity sudah
dipaksa secara struktural melalui operator $\text{Top}_{k}(\cdot)$.

\subsection{Pemetaan Frekuensi Karakteristik Bantalan (BPFx)}
\label{subsec:bab4_bpfx}

Setelah SAE dilatih, setiap fitur laten $f_j$ dievaluasi korelasinya
dengan amplitudo \emph{Hilbert envelope spectrum} sinyal getaran mentah
pada pita frekuensi $\pm 2$~Hz di sekitar frekuensi karakteristik
bantalan.
Empat frekuensi karakteristik yang digunakan:

\begin{equation}\label{eq:bab4_bpfo}
\text{BPFO} = \frac{n}{2} f_{r} \left(1 - \frac{d}{D}\cos\theta\right),
\end{equation}

\begin{equation}\label{eq:bab4_bpfi}
\text{BPFI} = \frac{n}{2} f_{r} \left(1 + \frac{d}{D}\cos\theta\right),
\end{equation}

\begin{equation}\label{eq:bab4_bsf}
\text{BSF} = \frac{D}{2d} f_{r} \left(1 - \left(\frac{d}{D}\cos\theta\right)^{2}\right),
\end{equation}

\begin{equation}\label{eq:bab4_ftf}
\text{FTF} = \frac{f_{r}}{2} \left(1 - \frac{d}{D}\cos\theta\right),
\end{equation}

\noindent di mana $n$ adalah jumlah elemen gelinding, $d$ adalah
diameter elemen, $D$ adalah \emph{pitch diameter}, $\theta$ adalah
sudut kontak (diasumsikan 0\textdegree\ untuk \emph{deep groove ball
bearing} tanpa beban aksial), dan $f_{r}$ adalah frekuensi rotasi
poros (Hz).
Nilai teoritis frekuensi karakteristik untuk kedua dataset ditampilkan
pada Tabel~\ref{tab:bab4_bpfx}.

\begin{table}[ht]
\centering
\caption{Frekuensi karakteristik teoritis untuk kedua dataset
(asumsi $\theta = 0$\textdegree).}
\label{tab:bab4_bpfx}
\footnotesize
\setlength{\tabcolsep}{4pt}
\begin{tabular}{@{}l l c c c c r r r@{}}
\toprule
\textbf{Dataset} & \textbf{Bantalan} & $n$ & $d$ (mm) & $D$ (mm) & $f_r$ (Hz)
  & \textbf{BPFO (Hz)} & \textbf{BPFI (Hz)} & \textbf{BSF (Hz)} \\
\midrule
PHM2012  & NSK 6804   & 13 & 3{,}5  & 25{,}5  & 30{,}0 & 168{,}24 & 221{,}76 & 107{,}23 \\
XJTU-SY  & LDK UER204 &  8 & 7{,}92 & 34{,}55 & 35{,}0 & 107{,}91 & 172{,}09 &  72{,}33 \\
\bottomrule
\end{tabular}
\end{table}

Korelasi Pearson dihitung antara vektor aktivasi fitur SAE dan amplitudo
\emph{envelope spectrum} menggunakan $n_{\text{corr}} = 300$ rekaman
dari \emph{training set}:

\begin{equation}\label{eq:bab4_pearson}
r_{j,\nu} = \frac{\sum_{i=1}^{n_{\text{corr}}} (f_{j,i} - \bar{f}_j)(A_{\nu,i} - \bar{A}_\nu)}
{\sqrt{\sum_{i}(f_{j,i}-\bar{f}_j)^2} \cdot \sqrt{\sum_{i}(A_{\nu,i}-\bar{A}_\nu)^2}},
\end{equation}

\noindent di mana $f_{j,i}$ adalah aktivasi fitur SAE ke-$j$ pada
rekaman ke-$i$, $A_{\nu,i}$ adalah amplitudo \emph{envelope spectrum}
pada pita frekuensi $\nu$ untuk rekaman ke-$i$, dan $\bar{(\cdot)}$
menyatakan nilai rata-rata.
Fitur SAE dianggap \emph{terpetakan} ke frekuensi $\nu$ bila
$|r_{j,\nu}| \geq 0{,}30$ (\emph{moderate correlation threshold}
berdasarkan literatur).

\section{Protokol Pelatihan}
\label{sec:bab4_training}

Semua model menggunakan konfigurasi pelatihan yang identik
(berkas \texttt{configs/\allowbreak train/\allowbreak cloud\_full\_75.yaml})
untuk memastikan komparabilitas yang adil.

\subsection{Pengaturan Inti}
\label{subsec:bab4_hyperp_training}

Tabel~\ref{tab:bab4_training} merangkum hiperparameter pelatihan utama.

\begin{table}[ht]
\centering
\caption{Konfigurasi pelatihan \texttt{cloud\_full\_75.yaml}.}
\label{tab:bab4_training}
\small
\begin{tabular}{@{}ll@{}}
\toprule
\textbf{Hiperparameter} & \textbf{Nilai} \\
\midrule
\texttt{max\_epochs}            & 75 \\
Optimizer                       & AdamW \\
\texttt{lr}                     & $8{,}0 \times 10^{-4}$ \\
\texttt{weight\_decay}          & $3{,}0 \times 10^{-4}$ \\
\texttt{xlstm\_lr\_mult}        & 0{,}75 (LR modul xLSTM $= 0{,}75 \times$ \texttt{lr}) \\
\texttt{freeze\_xlstm\_epochs}  & 5 (pembekuan xLSTM pada 5 epoch pertama) \\
\texttt{scheduler}              & \emph{Cosine annealing}, $T_{\text{max}} = 75$ \\
\texttt{warmup\_epochs}         & 5 \\
\texttt{monotonicity\_weight}   & 0{,}05 \\
\texttt{gradient\_clip\_val}    & 1{,}0 \\
\texttt{early\_stopping\_patience} & 9999 (dinonaktifkan) \\
\texttt{precision}              & \texttt{bf16-mixed} (Ampere+) \\
\bottomrule
\end{tabular}
\end{table}

\emph{Early stopping} dinonaktifkan agar semua model berjalan dengan
jumlah \emph{epoch} identik dan kurva pelatihan dapat dibandingkan
secara adil tanpa pemotongan asimetris.
\emph{Checkpoint} terbaik tetap disimpan berdasarkan \texttt{val/rmse}
minimum, sehingga evaluasi menggunakan model dengan performa validasi
optimal.

\subsection{Infrastruktur Komputasi}
\label{subsec:bab4_infra}

Konfigurasi perangkat keras dan \emph{environment} perangkat lunak yang
digunakan dalam seluruh eksperimen dirinci pada Tabel~\ref{tab:bab4_infra}.
Seluruh pelatihan dilakukan pada satu GPU (single-card) untuk menjamin
reproduktibilitas hasil antar-\emph{seed}.

\begin{table}[ht]
\centering
\caption{Infrastruktur pelatihan.}
\label{tab:bab4_infra}
\small
\begin{tabular}{@{}ll@{}}
\toprule
\textbf{Komponen} & \textbf{Spesifikasi} \\
\midrule
GPU             & NVIDIA A40 (46~GB VRAM) \\
\emph{Framework} & PyTorch Lightning 2.x \\
Presisi         & \texttt{bf16-mixed} (AMP) \\
\texttt{num\_workers} & 8 \\
Environment     & Python 3.11, \texttt{virtualenv} \\
\bottomrule
\end{tabular}
\end{table}

\section{Desain Pengulangan Multi-\emph{Seed}}
\label{sec:bab4_multiseed}

Untuk menjamin rigor statistik, setiap konfigurasi (model × dataset)
dijalankan dengan tiga \emph{seed} berbeda: 42, 43, dan 44.
Total 18 \emph{run} pelatihan dilakukan:

\begin{equation}\label{eq:bab4_totalruns}
\text{Total runs} = 3 \;\text{model} \times 2 \;\text{dataset} \times 3 \;\text{seed} = 18.
\end{equation}

Setiap \emph{run} menghasilkan satu \emph{best checkpoint} berdasarkan
\texttt{val/rmse} minimum.
Tabel final melaporkan rata-rata dan deviasi standar populasi atas
tiga \emph{seed}:

\begin{equation}\label{eq:bab4_mean_std}
\overline{M} \pm \sigma_{M} = \frac{1}{3}\sum_{s}\!M_{s}
  \;\pm\; \sqrt{\frac{1}{3}\sum_{s}\!\left(M_{s} - \overline{M}\right)^{2}}.
\end{equation}

Tiga \emph{seed} merupakan minimum yang direkomendasikan oleh praktik
terbaik literatur prognostik dan standar disertasi ITB untuk
pelaporan dengan format mean~$\pm$~std yang bermakna.

\section{Metrik Evaluasi}
\label{sec:bab4_metrik}

Empat metrik dihitung pada \emph{test set} untuk gambaran performa
yang multidimensional.

\subsection{RMSE dan MAE}
\label{subsec:bab4_rmse_mae}

\emph{Root Mean Squared Error} (RMSE) digunakan sebagai metrik utama:

\begin{equation}\label{eq:bab4_rmse}
\text{RMSE} = \sqrt{\frac{1}{N_{\text{test}}}\sum_{i=1}^{N_{\text{test}}} \left(\hat{r}_{i} - r_{i}\right)^{2}},
\end{equation}

\noindent dan \emph{Mean Absolute Error} (MAE) sebagai metrik pelengkap
yang lebih \emph{robust} terhadap \emph{outlier}:

\begin{equation}\label{eq:bab4_mae}
\text{MAE} = \frac{1}{N_{\text{test}}}\sum_{i=1}^{N_{\text{test}}} \left|\hat{r}_{i} - r_{i}\right|.
\end{equation}

\subsection{Koefisien Determinasi}
\label{subsec:bab4_r2}

\begin{equation}\label{eq:bab4_r2}
R^{2} = 1 - \frac{\text{SS}_{\text{res}}}{\text{SS}_{\text{tot}}}
      = 1 - \frac{\sum_{i}(\hat{r}_i - r_i)^2}{\sum_{i}(r_i - \bar{r})^2}.
\end{equation}

Nilai $R^{2} = 1$ menyatakan prediksi sempurna; nilai mendekati 0 setara
dengan prediktor \emph{mean}; nilai negatif menunjukkan performa lebih
buruk daripada prediktor \emph{mean}.
R\textsuperscript{2} negatif pada XJTU-SY disebabkan oleh \emph{test
split} yang sangat kecil (2 bantalan) yang menghasilkan variansi
SS\textsubscript{tot} kecil dan merupakan artefak statistik, bukan
indikasi performa buruk secara substantif.

\subsection{PHM Score}
\label{subsec:bab4_phmscore}

PHM Score merupakan \emph{scoring function} asimetris kompetisi
PHM2012 \citetitb{Nectoux2012PRONOSTIA}:

\begin{equation}\label{eq:bab4_phmscore}
S = \frac{1}{N_{\text{test}}} \sum_{i=1}^{N_{\text{test}}}
  \begin{cases}
  \exp\!\left(-\dfrac{\hat{r}_i - r_i}{\alpha_{1}}\right) & \text{bila } \hat{r}_i - r_i \leq 0 \\[6pt]
  \exp\!\left(-\dfrac{\hat{r}_i - r_i}{\alpha_{2}}\right) & \text{bila } \hat{r}_i - r_i > 0,
  \end{cases}
\end{equation}

\noindent dengan konstanta asimetri $\alpha_{1}$ dan $\alpha_{2}$
di mana $\alpha_{1} < \alpha_{2}$, sehingga prediksi terlalu cepat
(RUL yang terlalu besar/optimistis) dihukum lebih berat daripada
prediksi terlambat (RUL yang terlalu kecil/pesimistis).
Nilai PHM Score yang \emph{lebih tinggi} lebih baik ($S \in [0, 1]$).

\section{Rancangan Eksperimen}
\label{sec:bab4_eksperimen}

Eksperimen disusun dalam empat tahap yang saling membangun, mengikuti
struktur E1--E4 sesuai panduan metodologi.

\subsection{E1 — Validasi Konvergensi (Tahap 1)}
\label{subsec:bab4_e1}

\textbf{Tujuan:} Memastikan konvergensi model dan memvalidasi anggaran
\emph{epoch} yang cukup.

\textbf{Prosedur:} Mamba-xLSTM-Net dijalankan pada PHM2012 selama
200~\emph{epoch} (\emph{seed} 42) untuk mengidentifikasi \emph{best
epoch}.
Hasil: \emph{best checkpoint} pada \emph{epoch} 55 dengan
\texttt{val/RMSE = 0{,}113}; \emph{training loss} jenuh sekitar
\emph{epoch} 30; \emph{plateau val/RMSE} setelah \emph{epoch} 60.
\textbf{Kesimpulan:} Anggaran 75~\emph{epoch} pada Tahap 2 sudah
mencukupi untuk konvergensi optimal.

\subsection{E2 — Perbandingan Arsitektur Multi-\emph{Seed} (Tahap 2)}
\label{subsec:bab4_e2}

\textbf{Tujuan:} Mengevaluasi kinerja komparatif tiga model Tier-S
pada kedua dataset dengan tiga \emph{seed} berbeda.

\textbf{Prosedur:} Tiga model (Mamba-xLSTM-Net, N-BEATS-xLSTM-RUL,
SparseGate-TCN-RUL) × dua dataset (PHM2012, XJTU-SY) × 75~\emph{epoch}
× tiga \emph{seed} (42, 43, 44) = 18 \emph{run}.
Mamba-xLSTM-Net dan N-BEATS-xLSTM-RUL dijalankan \emph{paralel}
dalam satu GPU (batch 256 masing-masing); SparseGate-TCN-RUL
dijalankan \emph{solo} setelahnya (batch 512).
\emph{Wall-clock} per \emph{seed}: $\approx$85 menit.

\textbf{Keluaran:} Tabel mean~$\pm$~std per metrik per model per
dataset; analisis per-\emph{seed} dan \emph{best epoch}.

\subsection{E3 — Interpretabilitas SAE (Tahap 3)}
\label{subsec:bab4_e3}

\textbf{Tujuan:} Menganalisis struktur representasi laten
Mamba-xLSTM-Net melalui \emph{Top-$k$ Sparse Autoencoder}.

\textbf{Prosedur:} SAE (dimensi 1024, $k = 51$) dilatih selama
50~\emph{epoch} pada 20.000 \emph{hidden state} dari \emph{best
checkpoint} Mamba-xLSTM-Net per dataset (PHM2012: \emph{seed}~42,
\emph{epoch}~71; XJTU-SY: \emph{seed}~44, \emph{epoch}~70).
Analisis atribusi SHAP dan IG dilakukan pada fitur input.
Proyeksi UMAP dari ruang laten SAE divisualisasikan untuk mengidentifikasi
klaster semantik.

\textbf{Keluaran:} \emph{Checkpoint} SAE, kurva \emph{loss}, UMAP
klaster, SHAP global, IG per \emph{window}.

\subsection{E4 — Pemetaan BPFx (Tahap 4)}
\label{subsec:bab4_e4}

\textbf{Tujuan:} Menguji hipotesis H2 — apakah fitur SAE berkorelasi
dengan frekuensi karakteristik fisik bantalan.

\textbf{Prosedur:} Korelasi Pearson dihitung antara vektor aktivasi
tiap fitur SAE (1024 fitur) dan amplitudo \emph{Hilbert envelope
spectrum} pada pita $\pm 2$~Hz di sekitar BPFO, BPFI, BSF, FTF,
menggunakan $n = 300$ rekaman dari \emph{training set}.
Fitur dianggap \emph{terpetakan} bila $|r| \geq 0{,}30$.
\emph{Hit-rate} dihitung sebagai persentase fitur SAE yang terpetakan
dari total 1024 fitur.

\textbf{Keluaran:} Tabel \emph{hit-rate} per BPFx per dataset,
\emph{bar chart} \emph{hit-rate}, \emph{heatmap} korelasi top-50 fitur aktif.

% -----------------------------------------------------------------
\section{Studi Kasus Industri: PT~SKF~Indonesia (Pilar~0)}
\label{sec:bab4_studi_kasus_industri}
% -----------------------------------------------------------------

Untuk memvalidasi relevansi disertasi pada konteks produksi industri
nyata, dilakukan studi kasus pada lini produksi PT~SKF~Indonesia
(Pilar~0). Pilar~0 menempatkan kerangka kerja interpretabilitas yang
dikembangkan pada Pilar~1 dan Pilar~2 ke dalam tata-kerja
\emph{condition-based predictive maintenance} pabrik nyata, dengan
keluaran berupa klasifikasi tiga kelas kondisi (\emph{Normal},
\emph{Warning}, \emph{Critical}) yang langsung dapat ditindaklanjuti
oleh tim pemeliharaan.

\subsection{Konteks Industri dan Objek Studi}
\label{subsec:bab4_industri_konteks}

PT~SKF~Indonesia memproduksi bantalan gelinding untuk sektor otomotif
melalui beberapa tahapan permesinan: \emph{forging}, \emph{turning},
\emph{outer ring grinding}, \emph{inner ring grinding}, dan
\emph{honing} (lihat Gambar~\ref{fig:bab1_objek_industri}). Studi kasus
ini fokus pada dua mesin \emph{outer ring grinding} (OR1 dan OR2)
yang memproduksi cincin luar bantalan. Permasalahan utama adalah
tingginya \emph{scrap rate} akibat keausan \emph{cutting tool}
spindel yang menyebabkan parameter kualitas (\emph{vibration checking}
dan \emph{radial clearance checking}) keluar dari rentang spesifikasi.
Pendekatan pemeliharaan korektif konvensional (mengganti komponen
setelah terjadi penolakan produk) berdampak \emph{unplanned downtime}
dan kerugian material; klasifikasi kondisi tiga kelas berbasis sinyal
getaran spindel diharapkan memberikan peringatan dini sebelum produk
gagal.

\subsection{Akuisisi Data Multi-Sumber}
\label{subsec:bab4_akuisisi_data}

Sinyal getaran spindel diakuisisi oleh modul \emph{wireless} SKF~IMx-8
dengan akselerometer \emph{piezoelectric} (frekuensi sampling
$25{,}6$~kHz, \emph{snapshot} $1{,}28$~detik per perekaman, interval
$1$~menit). Data diteruskan ke SKF AWS Cloud melalui
\emph{gateway} pabrik. Pada paralel, data kualitas produk
(\emph{vibration checking} dan \emph{radial clearance checking}) di
\emph{end-of-line} stasiun dicatat oleh sistem MES pabrik dengan
\emph{timestamp} per produk. Integrasi dilakukan dengan kerangka
\emph{big data} \citetitb{Yan2017BigData}: \emph{event timestamp}
sinyal getaran disinkronisasi dengan \emph{batch timestamp} produk
hingga toleransi $\pm 30$~detik, lalu di-\emph{join} pada lapisan
\emph{data warehouse}. Skema akuisisi diilustrasikan oleh
Gambar~\ref{fig:bab4_pipeline_industri}.

\begin{figure}[ht]
  \centering
  \includegraphics[width=0.92\textwidth]{data_collection_pipeline.png}
  \caption{Pipa data \emph{end-to-end} studi kasus PT~SKF~Indonesia:
           sensor SKF~IMx-8 pada spindel mesin OR1/OR2, transmisi ke
           SKF AWS Cloud, integrasi dengan basis data kualitas produk
           pabrik, ekstraksi fitur, klasifikasi tiga kelas, dan
           penjelasan SHAP yang dikirim balik ke tim pemeliharaan.
           Sumber: PT~SKF~Indonesia~\citetitb{SKF2017TechReport}.}
  \label{fig:bab4_pipeline_industri}
\end{figure}

Setiap perekaman diolah menjadi vektor fitur 16 dimensi (9 fitur
domain waktu, 7 fitur domain frekuensi --- mengikuti rule domain
disertasi), dilengkapi tiga fitur kualitas turunan: \emph{vibration
checking max}, \emph{radial clearance mean}, dan \emph{rolling failure
ratio} dalam jendela 30 menit terakhir. Total dimensi fitur per
jendela adalah $19$.

\subsection{Skema Pelabelan Tiga Kelas}
\label{subsec:bab4_label_tiga_kelas}

Label tiga kelas dihitung secara otomatis menggunakan ambang
\emph{Progression Index} (PI) ternormalisasi yang sudah dijelaskan
pada Subbab~\ref{subsec:bab3_kls_label}. PI dihitung per perekaman
sebagai rasio kenaikan parameter \emph{velocity RMS} terhadap
\emph{baseline} bantalan baru, dikoreksi oleh kemiringan tren
\emph{moving window} 1~jam. Ambang $\text{PI} \in [0,\,0{,}20)$
menghasilkan label \emph{Normal}, $[0{,}20,\,0{,}50)$ menghasilkan
\emph{Warning}, dan $\geq 0{,}50$ menghasilkan \emph{Critical}.
Konsistensi label divalidasi terhadap ambang ISO~10816/20816 dan
\emph{cross-check} oleh teknisi senior pemeliharaan SKF.

\subsection{Pembagian Data dan Validasi Lintas Mesin}
\label{subsec:bab4_split_industri}

Pembagian data mengikuti dua skenario:

\begin{enumerate}
  \item \textbf{Skenario A (\emph{within-machine})}: data setiap mesin
        (OR1 atau OR2) dibagi 70/15/15 secara temporal --- 70\% rentang
        waktu paling awal sebagai \emph{train}, 15\% berikutnya sebagai
        \emph{validation}, dan 15\% paling baru sebagai \emph{test}. Skenario
        ini mengukur kinerja model pada mesin yang sama yang dipakai
        untuk pelatihan.
  \item \textbf{Skenario B (\emph{cross-machine})}: pelatihan dilakukan
        sepenuhnya pada satu mesin (mis. OR1) dan pengujian pada mesin
        lain (OR2), dan sebaliknya. Skenario ini mengukur generalisasi
        model lintas variasi unit.
\end{enumerate}

Statistik fitur (mean, std) di-\emph{fit} hanya pada \emph{training set}
masing-masing skenario, sesuai aturan domain. Skenario B berperan
penting untuk membuktikan bahwa keputusan klasifikasi tidak
\emph{overfit} pada karakteristik spesifik satu mesin.

\subsection{Benchmarking Enam Algoritma}
\label{subsec:bab4_benchmarking}

Enam algoritma kandidat (Subbab~\ref{subsec:bab3_kls_kandidat})
dilatih dengan protokol identik untuk memastikan perbandingan adil:

\begin{itemize}
  \item \textbf{Optimizer}: AdamW, \emph{learning rate} awal $10^{-3}$
        dengan \emph{cosine annealing} hingga $10^{-5}$.
  \item \textbf{Batch size}: $128$ jendela, $50$ \emph{epoch} maksimal,
        \emph{early stopping} dengan kesabaran $10$ \emph{epoch} pada
        $F1_{\text{macro}}$ validasi.
  \item \textbf{Loss}: \emph{class-weighted cross-entropy}
        (Persamaan~\ref{eq:bab3_xentropy}).
  \item \textbf{Multi-\emph{seed}}: tiga \emph{seed} ($42, 43, 44$);
        rerata $\pm$ \emph{standard deviation} dilaporkan.
\end{itemize}

Detail \emph{hyperparameter} per arsitektur dirinci pada
Lampiran~\ref{lamp:klasifikasi_industri}.

\subsection{Modul Interpretabilitas SHAP}
\label{subsec:bab4_shap_industri}

Untuk algoritma pemenang dari benchmarking, dihitung peringkat
\emph{global feature importance} via SHAP (Subbab~\ref{subsec:bab3_kls_shap}).
\emph{TreeExplainer} dipakai jika pemenang berbasis pohon (XGBoost),
sedangkan \emph{DeepExplainer} dipakai jika pemenang berbasis jaringan
saraf. Output utama adalah \emph{summary plot} per kelas dan
\emph{waterfall plot} untuk minimal lima sampel misklasifikasi
\emph{Critical} sebagai bahan \emph{root-cause analysis} bersama tim
pemeliharaan.

\subsection{Metrik Operasional Tambahan}
\label{subsec:bab4_metrik_operasional}

Selain $F1_{\text{macro}}$, akurasi, dan \emph{recall} per kelas,
dilaporkan dua metrik operasional yang menjadi syarat \emph{deployment}
di \emph{edge}: (i) \emph{inference latency} per jendela diukur pada
SKF~IMx-8 (target $\leq 10$~ms) dan (ii) jumlah parameter terlatih
sebagai proksi \emph{memory footprint}. Algoritma yang melampaui
$F1_{\text{macro}} = 0{,}85$ \emph{dan} memenuhi keduanya dianggap
kandidat siap-\emph{deployment}.
